\section{The quantitative theorem and normalized pairings}
\label{sec:background}

Singular integral estimates can detect the geometry of a measure even
when its support has no initial parametrization. The Riesz transform is
a particularly natural test: its kernel records the direction between
two points as well as their separation. We prove that, for an
integer-dimensional Ahlfors--David regular measure in the intermediate
codimensions, its $L^2$ boundedness forces a quantitatively rectifiable
support.

Throughout the paper, $d$ and $n$ are positive integers with $n\le d$,
and balls are open Euclidean balls. A Radon measure is
a Borel measure that is finite on compact sets and regular. For such a
measure $\mu$ on $\R^d$, write $E=\supp\mu$ and $D_E=\diam E$, allowing
$D_E=\infty$. When $E$ is empty, set $D_E=0$.

\begin{definition}[Ahlfors--David regularity]\label{def:ad}
A Radon measure $\mu$ on $\R^d$ is \emph{$n$-Ahlfors--David regular}
with constant $C_{\rm AD}\ge1$ if
\[
 C_{\rm AD}^{-1}r^n\le\mu(B(x,r))\le C_{\rm AD}r^n
 \qquad(x\in E,\quad 0<r\le D_E).
\]
We call these positive radii \emph{admissible}. The inequalities have
no instances when $D_E=0$.
\end{definition}

For $z\ne0$ let $K_n(z)=z/|z|^{n+1}$. The hard truncations of the
$n$-dimensional Riesz transform are
\begin{equation}\label{eq:hard-truncation}
 R_{\mu,\varepsilon}f(x)
 =\int_{|x-y|>\varepsilon} K_n(x-y)f(y)\,d\mu(y),\qquad\varepsilon>0.
\end{equation}
We say that the transform is bounded on $L^2(\mu)$ with bound
$C_{\rm R}$ if, for every real scalar $f\in L^2(\mu)$ and every
$\varepsilon>0$, the integral in \eqref{eq:hard-truncation} defines an
$\R^d$-valued $L^2(\mu)$ function and
\begin{equation}\label{eq:riesz-bound}
 \|R_{\mu,\varepsilon}f\|_{L^2(\mu;\R^d)}
 \le C_{\rm R}\|f\|_{L^2(\mu)}.
\end{equation}
No pointwise principal-value hypothesis is imposed. For a nontrivial
regular support, the integrals in \eqref{eq:hard-truncation} are
absolutely convergent for every $x$ and every $f\in L^2(\mu)$;
this follows from Lemma~\ref{lem:global-growth} below.

\begin{definition}[Uniform rectifiability]\label{def:ur}
An $n$-Ahlfors--David regular measure $\mu$ on $\R^d$ is
\emph{uniformly $n$-rectifiable} if there are $\theta>0$ and $M<\infty$
such that, for every $x\in E$ and $0<r\le D_E$, there is an
$M$-Lipschitz map
\[
 g:B_{\R^n}(0,r)\longrightarrow\R^d,
 \qquad
 \mu\bigl(B(x,r)\cap g(B_{\R^n}(0,r))\bigr)\ge\theta r^n.
\]
The constants are chosen before $x$ and $r$.
\end{definition}

\begin{theorem}\label{thm:main}
Let $d\ge4$ and $2\le n\le d-2$ be integers. Let $\mu$ be an
$n$-Ahlfors--David regular Radon measure on $\R^d$ satisfying
\eqref{eq:riesz-bound}. Then $\mu$ is uniformly $n$-rectifiable in the
sense of Definition~\ref{def:ur}. Its constants $\theta$ and $M$ can
be chosen in terms of $d,n,C_{\rm AD}$ and $C_{\rm R}$ alone.
\end{theorem}

\subsection{History and the analytic question}

David and Semmes developed uniform rectifiability as a quantitative
form of rectifiability suited to singular integrals
\cite{DavidSemmes1991,DavidSemmes1993}. Among its equivalent
descriptions are big pieces of Lipschitz images and a packing condition
on balls where the support fails to approximate a plane from both
sides. Their question asks whether the boundedness of the Riesz
transform alone forces this geometry for an integer-dimensional
regular measure.

Mattila, Melnikov and Verdera established the planar implication using
the relation between the Cauchy integral, curvature and uniform
rectifiability \cite{MattilaMelnikovVerdera1996}. Curvature methods
also give the one-dimensional Riesz-transform implication in higher
ambient dimension; see \cite[Theorem~4]{Farag2000}.
Nazarov, Tolsa and Volberg proved the codimension-one implication
\cite{NazarovTolsaVolberg2014}.
They identify reliance on a maximum principle, for which they report
no known higher-codimension analogue, as the main obstacle to extending
their argument \cite[Section~1]{NazarovTolsaVolberg2014}.
Mas and Tolsa obtained a characterization in every integer dimension
$1\le n<d$ by an $L^2$ bound for the $\rho$-variation of the Riesz
truncations, with $\rho>2$ \cite[Theorem~1.3]{MasTolsa2014Variation}.
That hypothesis bounds the pointwise supremum, over decreasing positive
sequences $(\varepsilon_j)$, of
\[
 \left(\sum_j
   |R_{\mu,\varepsilon_{j+1}}f-R_{\mu,\varepsilon_j}f|^\rho
       \right)^{1/\rho}
\]
in $L^2(\mu)$. It controls changes across scales, rather than each
truncation separately. The present hypothesis is only the common bound
\eqref{eq:riesz-bound}; no variation estimate is assumed.
The remaining range for this bound-alone question,
$1<n<d-1$, is stated explicitly in \cite[Question~4.4]{Tolsa2026ICM}.
Theorem~\ref{thm:main} gives a positive answer in that range, with the
hard-truncation convention and the ball parametrizations stated
above.

Principal-value criteria address a different analytic hypothesis.
For a set of finite $\H^n$ measure, Tolsa characterized ordinary
rectifiability by almost-everywhere existence of Riesz principal values
\cite[Theorem~1.1]{Tolsa2008PrincipalValues}. That result neither assumes
only \eqref{eq:riesz-bound} nor asserts the quantitative ball-map
conclusion considered here.

Two parts of the quantitative-rectifiability literature are especially
relevant to the proof. The bilateral weak geometric lemma of David and
Semmes converts geometric packing into Lipschitz images; a precise
Euclidean formulation is recalled in
\cite[Theorem~2.5]{Tolsa2015Uniform}. Reflectionless-measure methods
organize compactness limits for which the Riesz transform has zero
oscillation against mean-zero tests; see
\cite{JayeNazarov2018I}. We prove the particular limiting and rigidity
statements needed here, including their tail normalizations.

Flat subballs and stability of flatness also form the geometric route
in Tolsa's theorem for uniform measures
\cite[Section~3]{Tolsa2015Uniform}. There the ball masses are exactly
proportional to $r^n$, and the stability step uses Preiss's rigidity at
infinity, recalled in \cite[Theorem~2.4]{Tolsa2015Uniform}.
Here AD regularity gives two-sided comparisons rather than exact ball
masses. The normalized-height energy and rigidity argument below
provides the required propagation without a uniform-measure rigidity
input.

Section~\ref{sec:consequences} applies the uniform-rectifiability
conclusion to obtain variation bounds and almost-everywhere principal
values for the Riesz transform and the odd $C^2$ kernels specified there.

\subsection{The main mechanism}

The analytic difficulty is to propagate geometric flatness to smaller
scales. We measure the error by the \emph{excess}, the scale-normalized
$L^2$ distance from an affine $n$-plane. Normal heights are the coordinates
of the displacement perpendicular to that plane. Weak convergence to a
plane makes the excess tend to zero; propagation requires an error small
relative to the initial excess. We therefore divide the normal heights
by that initial scale before passing to a limit.
This second limit is taken only when the Riesz pairings against compact
mean-zero Lipschitz tests are sufficiently small relative to the height
scale to vanish after normalization. Here the mean is taken against
the current measure, not against a presumed limiting plane.
The use of normal components to measure geometric deviation has a
related quantitative form on small-slope Lipschitz graphs in
\cite[Theorem~1.3]{Tolsa2008PrincipalValues}. Our compactness argument
instead works on moving AD-regular measures before any graph
parametrization has been established.

Testing the transform against an unnormalized normal height times a
cutoff produces a positive fractional energy. This estimate gives
compactness of the normalized heights through their second moments
and their integrals against Lipschitz tests. It also controls the
singular diagonal when passing to a limiting equation. The equation
holds against mean-zero tests and is therefore interpreted modulo
constants. Its order-one Fourier symbol, together with weighted
integrability at infinity, forces each limiting height to be affine. Approximation
by an affine graph now improves the original excess at the required
relative scale.

The remaining argument counts failures of this propagation on dyadic
cells of the support, nested pieces at successively halved scales.
The operator bound supplies descendants whose support is bilaterally
close to an $n$-plane, except for a family with a uniform packing bound.
Starting from such a flat descendant, a later failure of flatness forces
an intervening cell with detectable transform oscillation. A finite
counting argument controls these failures even when their scales are
widely separated. The David--Semmes theorem then gives the maps in
Definition~\ref{def:ur}.

The proof also gives a compactness principle for heights with bounded
local fractional energy on measures approaching plane measure. Combined
with the limiting rigidity equation, it yields affine approximation
while both the measure and its supporting set vary.
The renormalized equation is proved with the precise exterior
integrability that admits affine functions. This is the relevant
polynomial-growth setting for fractional operators; compare
\cite{DipierroSavinValdinoci2019}.

The rest of this section fixes the growth, pairing and dyadic conventions
and gives a precise proof map. Section~\ref{sec:flat-balls} proves the two packing estimates.
Section~\ref{sec:planar-limits} identifies the limiting plane measure.
Sections~\ref{sec:height-compactness} and~\ref{sec:fractional-rigidity}
establish compactness and affine rigidity of the normalized heights.
Section~\ref{sec:excess-propagation} proves uniform excess propagation,
and Section~\ref{sec:packing-endpoint} completes the packing argument
and Theorem~\ref{thm:main}. Section~\ref{sec:consequences} derives the
variation and principal-value consequences. Appendix~\ref{app:geometric-conventions}
records the geometric convention changes used at the endpoint.
